% !TeX root = /home/gsonoda/project/unicamp/unicamp_mo443_mc920/docs/relatorio.tex
\documentclass[12pt,a4paper]{article}

\usepackage[T1]{fontenc}
\usepackage[utf8]{inputenc}
\usepackage[brazilian]{babel}
\usepackage{lmodern}
\usepackage[a4paper,margin=2.5cm]{geometry}
\usepackage{amsmath}
\usepackage{amssymb}
\usepackage{graphicx}
\usepackage{float}
\usepackage{hyperref}
\usepackage{xcolor}

\hypersetup{
  colorlinks=true,
  linkcolor=blue,
  urlcolor=blue,
  pdftitle={Relatorio do Trabalho 1},
  pdfauthor={Gustavo Shoiti Sonoda}
}

\renewcommand{\contentsname}{Sumário}
\setlength{\parindent}{0pt}
\setlength{\parskip}{0.7em}
\linespread{1.05}

\newcommand{\placeholder}[1]{\textcolor{red}{\textbf{[Preencher: #1]}}}
\newcommand{\arquivo}[1]{\texttt{#1}}
\newcommand{\campo}[2]{\textbf{#1.} #2\par}

\begin{document}
\hypersetup{pageanchor=false}
\pagenumbering{gobble}

\begin{titlepage}
  \centering

  {\Large \textbf{Relatório de Implementação do Trabalho 1}\par}
  \vspace{0.5cm}
  {\large Introdução ao Processamento Digital de Imagens (MC920/MO443)\par}
  \vspace{1.2cm}

  {\large Relatório de implementação desenvolvido para a disciplina\\
  Introdução ao Processamento Digital de Imagens.\par}

  \vspace{1.5cm}

  \begin{flushleft}
    \textbf{Professor:} Hélio Pedrini\\
    \textbf{Aluno:} Gustavo Shoiti Sonoda\\
    \textbf{RA:} \placeholder{informar o RA}
  \end{flushleft}

  \vfill

  Campinas -- São Paulo\\
  \today
\end{titlepage}

\tableofcontents
\newpage
\hypersetup{pageanchor=true}
\pagenumbering{arabic}

\section{Introdução}

O objetivo deste trabalho é realizar processamentos básicos em imagens digitais, conforme especificado no enunciado do Trabalho 1 da disciplina MC920/MO443. As atividades propostas abrangem operações geométricas, transformações pontuais, manipulação de imagens coloridas, decomposição em planos de bits, combinação de imagens, quantização e filtragem espacial. Sempre que pertinente, a implementação deve explorar vetorização de operações, respeitando o formato de entrada e saída em imagens PNG e o ambiente de execução em Linux.

Este relatório foi organizado de modo semelhante ao modelo observado no relatório de referência, separando a descrição do ambiente e das decisões de implementação da análise dos resultados experimentais. A seção de materiais e métodos apresenta a organização do projeto, os dados de entrada e a estratégia geral adotada. Em seguida, a seção de resultados e discussão percorre cada item do enunciado, mantendo a numeração do trabalho para facilitar a conferência entre implementação, relatório e requisitos solicitados.

\section{Materiais e Métodos}

\subsection{Requerimentos}

De acordo com o enunciado, os programas serão executados em ambiente Linux e os formatos de entrada e saída devem ser rigorosamente respeitados. As imagens utilizadas no trabalho estão no formato PNG, tanto para leitura quanto para gravação dos resultados. Além disso, o enunciado destaca que, quando pertinente, a vetorização de comandos deve ser empregada nas operações.

\campo{Linguagem e bibliotecas}{\placeholder{Informar a linguagem escolhida, a versão utilizada e as bibliotecas empregadas na leitura, processamento e escrita das imagens.}}
\campo{Entrada de dados}{As imagens de entrada estão no formato PNG, conforme especificado no trabalho. \placeholder{Indicar quais imagens foram usadas nos experimentos e de onde elas foram obtidas.}}
\campo{Saída de dados}{As imagens geradas pelo programa também devem ser salvas em PNG. Resultados intermediários podem ser exibidos em tela ou armazenados em arquivos auxiliares.}

\subsection{Organização do projeto}

O projeto foi estruturado para separar a leitura e escrita das imagens das rotinas de processamento. Essa divisão facilita o reaproveitamento de código entre os itens do trabalho e torna mais clara a validação de cada transformação isoladamente. \placeholder{Listar aqui os arquivos reais do projeto e a responsabilidade de cada um, por exemplo: arquivo principal, módulo de operações em imagens monocromáticas, módulo de operações em imagens RGB, utilitários e pasta de resultados.}

\campo{Organização de arquivos}{\placeholder{Descrever como as pastas de entrada, saída e código foram organizadas no repositório.}}
\campo{Representação das imagens}{\placeholder{Explicar como as imagens são armazenadas em memória, indicando a convenção de índices, o tipo dos dados e a forma de acesso aos canais de cor.}}
\campo{Convenções adotadas}{\placeholder{Comentar tratamento de overflow, clipping para o intervalo $[0,255]$, arredondamento, normalização e política para tratamento de bordas.}}

\subsection{Estratégia geral de implementação}

De forma geral, cada item do trabalho segue um mesmo fluxo: leitura da imagem PNG, aplicação da transformação solicitada, validação do resultado e gravação da nova imagem em disco. Esse fluxo é especialmente útil neste trabalho porque várias operações atuam diretamente sobre os valores da matriz da imagem, enquanto outras dependem de relações espaciais entre pixels vizinhos ou de combinações lineares entre canais de cor.

\campo{Esboço do trabalho}{A implementação pode ser entendida em três blocos principais: (i) infraestrutura de leitura e escrita de imagens, (ii) implementação das operações pedidas no enunciado e (iii) geração das saídas utilizadas para comparação visual e discussão dos resultados.}
\campo{Estrutura do relatório}{O relatório final deve registrar, para cada item, o objetivo da operação, a estratégia de implementação, os parâmetros utilizados, os testes executados, os resultados obtidos e eventuais limitações observadas.}

\subsection{Procedimento de testes}

Os testes devem verificar tanto a corretude funcional das operações quanto a qualidade visual dos resultados. Em itens geométricos, a validação pode ser feita pela inspeção da disposição dos pixels e das dimensões da imagem. Em itens pontuais, a comparação entre valores de entrada e saída e a observação do histograma visual são úteis. Em filtragem e manipulação de cor, a análise visual das imagens geradas é essencial para discutir suavização, realce de contornos, preservação de detalhes e saturação.

\campo{Conjunto de testes}{\placeholder{Descrever quais imagens foram usadas em cada item e justificar a escolha delas.}}
\campo{Registro das saídas}{\placeholder{Indicar o padrão de nomes dos arquivos gerados e onde eles foram armazenados.}}
\campo{Critérios de validação}{\placeholder{Explicar como cada transformação foi validada: inspeção visual, conferência de dimensões, comparação de intensidade, testes com parâmetros diferentes, entre outros.}}

\section{Resultados e Discussão}

Esta seção apresenta os resultados obtidos em cada item do Trabalho 1. A organização abaixo segue a numeração do enunciado para manter o relatório alinhado com a especificação oficial da atividade.

\subsection{Item 1.1 -- Rotação de Imagens em Múltiplos de 90 Graus}

\campo{Objetivo}{Dada uma imagem monocromática, implementar a rotação nos ângulos de 90 graus no sentido horário, 180 graus e 270 graus no sentido horário, sem utilizar funções prontas de rotação das bibliotecas.}
\campo{Implementação}{\placeholder{Descrever o mapeamento direto dos índices da matriz para cada ângulo e explicar como largura e altura foram tratadas nas rotações.}}
\campo{Testes executados}{\placeholder{Comparar a imagem original com as três versões rotacionadas e registrar como a corretude foi conferida.}}
\campo{Discussão}{\placeholder{Comentar simplicidade da solução, custo computacional e cuidados tomados para evitar erros de indexação.}}

\subsection{Item 1.2 -- Ampliação de Imagem por Replicação}

\campo{Objetivo}{A partir de uma imagem monocromática, gerar versões ampliadas por replicação de pixels com fatores $2$ e $4$.}
\campo{Implementação}{\placeholder{Explicar como cada pixel original foi replicado para formar blocos na imagem ampliada.}}
\campo{Testes executados}{\placeholder{Descrever como foram conferidas as novas dimensões e a preservação do padrão visual da imagem.}}
\campo{Discussão}{\placeholder{Comentar o efeito serrilhado esperado da replicação e comparar com o que seria obtido por interpolação.}}

\subsection{Item 1.3 -- Esboço a Lápis}

\campo{Objetivo}{Implementar um efeito de esboço a lápis em uma imagem colorida por meio dos passos descritos no enunciado: conversão para níveis de cinza, aplicação de filtro gaussiano de desfoque e divisão da imagem em tons de cinza pela versão desfocada.}
\campo{Implementação}{\placeholder{Descrever as etapas utilizadas, indicando como a imagem foi convertida para cinza, como o desfoque gaussiano foi aplicado e como a divisão entre imagens foi tratada numericamente.}}
\campo{Parâmetros adotados}{O enunciado sugere o uso de uma máscara gaussiana de $21 \times 21$ pixels. \placeholder{Indicar se esse valor foi mantido ou se houve alguma adaptação.}}
\campo{Testes executados}{\placeholder{Informar quais imagens foram usadas para avaliar contornos, sombras e regiões com diferentes níveis de textura.}}
\campo{Discussão}{\placeholder{Comentar a aparência do efeito obtido e a influência do desfoque no realce das estruturas da imagem.}}

\subsection{Item 1.4 -- Ajuste de Brilho}

\campo{Objetivo}{Aplicar correção gama para ajustar o brilho de uma imagem monocromática $A$ de entrada e gerar uma imagem monocromática $B$ de saída, utilizando diferentes valores de $\gamma$.}
\campo{Implementação}{\placeholder{Descrever a transformação, incluindo a conversão do intervalo $[0,255]$ para $[0,1]$, a aplicação da equação e o retorno ao intervalo original.}}
\campo{Equação utilizada}{Uma forma de implementação sugerida pelo enunciado é:
\[
  B = 255 \cdot \left(\frac{A}{255}\right)^{1/\gamma}.
\]
\placeholder{Confirmar se essa foi exatamente a forma utilizada no código.}}
\campo{Testes executados}{\placeholder{Registrar os valores de $\gamma$ avaliados, por exemplo $1{,}5$, $2{,}5$ e $3{,}5$, e comparar visualmente os resultados.}}
\campo{Discussão}{\placeholder{Explicar como a variação de $\gamma$ afetou o brilho e o contraste da imagem.}}

\subsection{Item 1.5 -- Binarização por Limiar}

\campo{Objetivo}{Dada uma imagem monocromática $A$, gerar uma imagem binária $B$ a partir de um limiar fixo $T$.}
\campo{Implementação}{A conversão é dada por:
\[
  B =
  \begin{cases}
    255, & \text{se } A > T,\\
    0, & \text{caso contrário.}
  \end{cases}
\]
\placeholder{Explicar como o limiar foi parametrizado e como a operação foi aplicada à imagem.}}
\campo{Testes executados}{\placeholder{Informar os limiares testados e justificar a escolha do valor apresentado nos resultados.}}
\campo{Discussão}{\placeholder{Discutir perdas de detalhe, sensibilidade a ruído e separação entre objeto e fundo.}}

\subsection{Item 1.6 -- Mosaico}

\campo{Objetivo}{Construir um mosaico de $4 \times 4$ blocos a partir de uma imagem monocromática, seguindo a nova ordem de blocos mostrada no enunciado.}
\campo{Implementação}{\placeholder{Explicar como a imagem foi particionada em 16 blocos e como o remapeamento para a nova ordem foi realizado.}}
\campo{Testes executados}{\placeholder{Descrever como foi validada a correspondência entre a ordem original e a ordem final dos blocos.}}
\campo{Discussão}{\placeholder{Comentar o efeito visual obtido e os cuidados necessários quando a imagem precisa ser dividida exatamente em blocos de mesmo tamanho.}}

\subsection{Item 1.7 -- Alteração de Cores}

\campo{Objetivo}{Implementar um filtro para simular o efeito de fotografias antigas por meio da aplicação de uma transformação linear nas cores da imagem.}
\campo{Implementação}{Cada pixel deve ser multiplicado pela matriz:
\[
  \begin{bmatrix}
    0.393 & 0.769 & 0.189\\
    0.349 & 0.686 & 0.168\\
    0.272 & 0.534 & 0.131
  \end{bmatrix}.
\]
Após a transformação, os valores devem ser limitados ao intervalo de $0$ a $255$. \placeholder{Descrever como o cálculo foi implementado no código.}}
\campo{Testes executados}{\placeholder{Indicar quais imagens foram usadas para observar tons de pele, vegetação, contraste e saturação.}}
\campo{Discussão}{\placeholder{Comentar o efeito visual produzido pela transformação e o impacto do clipping em $255$.}}

\subsection{Item 1.8 -- Transformação de Imagens Coloridas}

\campo{Objetivo}{Aplicar as transformações pedidas no enunciado para imagens RGB, contemplando tanto a transformação componente a componente quanto a conversão para uma única banda por média ponderada.}
\campo{Implementação do item (a)}{A transformação componente a componente é dada por:
\[
  R' = 0.393R + 0.769G + 0.189B,
\]
\[
  G' = 0.349R + 0.686G + 0.168B,
\]
\[
  B' = 0.272R + 0.534G + 0.131B.
\]
Caso $R'$, $G'$ ou $B'$ ultrapassem $255$, seus valores devem ser limitados nesse máximo. \placeholder{Explicar como essa regra foi implementada.}}
\campo{Implementação do item (b)}{A imagem com uma única banda é obtida por:
\[
  I = 0.2989R + 0.5870G + 0.1140B.
\]
\placeholder{Descrever o formato da saída e como a conversão foi armazenada.}}
\campo{Testes executados}{\placeholder{Explicar como foram comparadas a saída colorida transformada e a versão em intensidade única.}}
\campo{Discussão}{\placeholder{Comentar diferenças visuais entre a transformação colorida e a imagem gerada pela média ponderada.}}

\subsection{Item 1.9 -- Planos de Bits}

\campo{Objetivo}{Extrair os planos de bits de uma imagem monocromática, com base na representação de cada pixel como polinômio de base 2.}
\campo{Fundamentação}{Para uma imagem monocromática com $m$ bits, os níveis de cinza podem ser representados por:
\[
  a_{m-1}2^{m-1} + a_{m-2}2^{m-2} + \ldots + a_1 2^1 + a_0 2^0.
\]
O plano de bits de ordem $0$ é formado pelos coeficientes $a_0$ e o plano de ordem $m-1$ pelos coeficientes $a_{m-1}$.}
\campo{Implementação}{\placeholder{Explicar como cada pixel foi decomposto em bits e como cada plano foi reconstruído em forma de imagem.}}
\campo{Testes executados}{\placeholder{Descrever a análise dos planos de bits $0$, $4$ e $7$, conforme ilustrado no enunciado.}}
\campo{Discussão}{\placeholder{Comentar a contribuição visual de cada plano de bits e a diferença entre planos menos e mais significativos.}}

\subsection{Item 1.10 -- Combinação de Imagens}

\campo{Objetivo}{Combinar duas imagens monocromáticas de mesmo tamanho por meio de média ponderada dos níveis de cinza.}
\campo{Implementação}{\placeholder{Registrar a expressão utilizada, por exemplo $C = \alpha A + (1-\alpha)B$, e explicar os pesos escolhidos.}}
\campo{Testes executados}{\placeholder{Descrever os pares de imagens usados e os resultados obtidos para pesos como $0{,}2/0{,}8$, $0{,}5/0{,}5$ e $0{,}8/0{,}2$.}}
\campo{Discussão}{\placeholder{Comentar mistura visual, preservação de detalhes e transição gradual entre as duas imagens.}}

\subsection{Item 1.11 -- Transformação de Intensidade}

\campo{Objetivo}{Dada uma imagem monocromática, realizar as transformações pedidas no enunciado. Isso inclui gerar o negativo da imagem, converter o intervalo de intensidades para o intervalo $[100,200]$, inverter os pixels das linhas pares, espelhar a metade superior na parte inferior e aplicar espelhamento vertical em toda a imagem.}
\campo{Implementação}{\placeholder{Descrever cada transformação separadamente, deixando claro como o mapeamento de intensidades e as operações geométricas foram realizadas.}}
\campo{Testes executados}{\placeholder{Explicar como cada subitem foi validado visualmente e, quando necessário, por conferência manual de pixels ou linhas.}}
\campo{Discussão}{\placeholder{Comparar os efeitos do negativo, do remapeamento de intensidades e dos diferentes espelhamentos.}}

\subsection{Item 1.12 -- Quantização de Imagens}

\campo{Objetivo}{Representar uma imagem monocromática com diferentes números de níveis de cinza, explorando a relação entre quantização e profundidade da imagem.}
\campo{Implementação}{\placeholder{Descrever o procedimento utilizado para gerar versões com 256, 64, 32, 16, 8, 4 e 2 níveis de cinza.}}
\campo{Testes executados}{\placeholder{Indicar como a perda progressiva de níveis foi analisada nas imagens quantizadas.}}
\campo{Discussão}{\placeholder{Comentar a degradação visual causada pela redução do número de níveis e em que ponto a perda de qualidade se torna mais evidente.}}

\subsection{Item 1.13 -- Filtragem de Imagens}

\campo{Objetivo}{Aplicar os filtros $h_1$ a $h_{11}$ em uma imagem digital monocromática e explicar os efeitos produzidos por cada máscara.}
\campo{Implementação}{No processo de filtragem, utiliza-se a operação de convolução entre a máscara e a imagem. O enunciado também solicita que os filtros $h_3$ e $h_4$ sejam aplicados individualmente e de forma combinada, por meio da expressão:
\[
  \sqrt{(h_3)^2 + (h_4)^2}.
\]
\placeholder{Explicar como a convolução foi implementada, qual tratamento de borda foi adotado e como as saídas foram normalizadas ou limitadas.}}
\campo{Testes executados}{\placeholder{Descrever quais imagens foram usadas para avaliar suavização, realce, detecção de bordas e outros efeitos produzidos pelos filtros.}}
\campo{Discussão}{\placeholder{Apresentar uma interpretação para cada máscara, indicando se o comportamento principal observado foi suavização, realce, aguçamento, detecção de bordas ou outro efeito.}}

\section{Conclusão}

\subsection{Limitações}

\placeholder{Listar as principais limitações da implementação, como tratamento de bordas, uso de parâmetros fixos, dependência do tamanho da imagem, ausência de otimizações específicas ou situações especiais não tratadas pelo programa.}

\subsection{Síntese final}

O Trabalho 1 reúne um conjunto bastante variado de operações de processamento digital de imagens, cobrindo desde transformações simples por pixel até operações espaciais baseadas em vizinhança. Essa variedade torna o trabalho particularmente útil para comparar como diferentes classes de transformação afetam a aparência final da imagem e quais cuidados de implementação surgem em cada caso.

\placeholder{Substituir este parágrafo por uma conclusão final baseada nos seus resultados, destacando os itens mais desafiadores, os principais aprendizados e as observações mais importantes obtidas nos testes.}

\end{document}
